\section{Implementation details}\label{app:impl}
% =============================================================================================
% Owner: implementation writer. Sources: src/synthval (twosample/nets.py, critic.py, c2st.py,
% mmd.py, signature.py, relative.py; models/univariate.py, multivariate.py; pipeline.py),
% experiments/*.py (FULL configs), experiments/README.md, run_full.sh, results/logs/*.log,
% results/*/meta*.json, docs/DESIGN.md (Section 4 and change log).
% Derived numbers (computed by scripts, not typed): early-stopping checkpoints from
% results/e2/e2_main_*_raw.csv and results/e3/e3_pooled_raw.csv; HAC lags and wild-bootstrap
% lengths from results/e3/e3_pooled_raw.csv; summed replication times and wall clock from
% results/logs (lanes A-D and e2_main.log, 15 November runs only); Merton truncation, Heston
% path count, signature dimensions and block sizes from the library functions.
% Architecture ablation: results/e5arch/meta*.json (h0_reps 200, sim_reps 50, power_reps 20),
% results/logs/e5arch.log (H0 extension 120 -> 180 blocks of 5; summed block times 0.865 + 3.21 +
% 3.01 + 10.77 + 6.45 = 24.3 h; wall 09:03:06-12:06:39 = 3.06 h), tab:e5-arch note (66 s vs 9 s).
% Provenance (app:impl-avail): git, code, script, data and packages fields of every
% results/{e1,e2,e2dep,e3,e4,e5,e5arch}/meta*.json; sha256sum of experiments/*.py and of the data
% file (all equal to the recorded hashes on 16 Nov 2026); .gitignore (results/*_quick/,
% results/logs/, checkpoints untracked); results/e3_quick/meta*.json and checkpoint file times
% (first pooled smoke run 14:17, first D2 smoke run 16:18 on 15 Nov); docs/DESIGN.md change log
% (learning-rate check, classifier pilot, architecture pilot); get_device in twosample/nets.py.
% =============================================================================================

This appendix records the implementation choices that \cref{sec:method,sec:data} leave out, at
the level needed to reproduce the reported numbers. It describes the code that produced them:
the Python package \texttt{synthval} and the experiment scripts of E1--E5. The default settings
of the learned tests are collected in \cref{tab:method-defaults} and the generators are defined
in \cref{sec:data-generators}. \Cref{app:impl-avail} states which code state produced each
result, what the run records contain, and what cannot be redistributed or reproduced from the
repository.

% ---------------------------------------------------------------------------------------------
\subsection{Networks, standardisation and training}\label{app:impl-nets}

\paragraph{Networks.}
The networks of \cref{sec:method-critic} (the perceptron for univariate windows, the
convolutional network with the 30 assets as input channels for the $(N,L)=(30,5)$ windows of
E4) are written in PyTorch \citep{paszke_2019_pytorch}, run in single precision on the CPU with
one thread per process, and start from PyTorch's default initialisation. The E4 network has
52{,}097 trainable parameters and is trained on 60 to 63 windows per sample (both counted by
\texttt{paper/scripts/e4\_reliability.py} from the network definition and the E4 results). The \WC{}, \LRC{} and
C2ST tests use the same networks and differ only in the loss. After training, the network is
frozen and evaluated without gradients, in batches of 8192 windows.

An ablation (script \texttt{e5\_architecture.py}) replaces the perceptron by a time-pooled
convolutional critic: two convolutions with 64 channels, kernel size 5 and ``same'' padding,
each followed by a LeakyReLU(0.2), the mean of the feature map over the $L$ positions, a dense
layer of width 64, a LeakyReLU and a linear output. Its output is an average of one learned
function of nine-day patches, so it uses every day of a window but cannot see where in the
window a pattern occurs. With the data and test randomness of the original runs, the ablation
reruns the \WC{} test in ten cells of E2 ($n=500$; 200 replications in each of the two null
cells and 50 in each of the eight alternative cells) and in two alternative cells of E2-dep
($n=2000$; 50 replications each), and it reruns the E3 protocol with E3's five seeds and designs
D1 and D2 with 20 replications of five seeds each, paired with the first 20 replications of E3.
The two null cells were first run with 50 replications, like the others, and then extended to
200. The results are in \cref{tab:e5-arch,tab:e5-arch-real} (\cref{app:robust}); a test with
the pooled critic takes on average 66\,s on one thread, against 9\,s with the perceptron. The
architecture was selected after the real-data results of the perceptron were known and after a
pilot on four E2 designs at $n=500$, which is recorded only in the change log of the protocol
and is not in the repository (\cref{app:impl-avail}). The simulated gains of the ablation are
therefore measured partly on the experiment and sample size used for the selection. The
ablation is post hoc: the runs of E1--E5 described in \cref{app:impl-compute}, and every verdict
of \cref{sec:emp}, use the perceptron for univariate windows.

\paragraph{Standardisation.}
The affine map $A$ of \cref{sec:method-critic} is estimated on the pooled training windows; in
the anchored test these are all $\Rnew$ windows and the training windows of the pool, without
weights. Means and standard deviations are taken per coordinate for univariate windows and per
asset channel, pooled over the $L$ positions, for multivariate windows. A standard deviation
below $10^{-12}$ is replaced by one.

\paragraph{Training of the critic.}
One forward pass per step evaluates the network on the $b=128$ real windows, the $b$ synthetic
windows and the $b$ interpolates $\tilde x_r$ of \eqref{eq:wgan-gp}. The gradient
$\nabla g(\tilde x_r)$ is obtained by automatic differentiation with respect to the
standardised input and kept in the computational graph, so that the penalty itself is
differentiated. The Adam step \citep{kingma_2015_adam} has no weight decay, there is no weight
or gradient clipping and no early stopping, so every critic runs the full 1000 steps. In the anchored test
the positive class ($\Rnew$) is sampled uniformly and the pool with the probabilities of
\cref{sec:method-relative}, by multinomial sampling with replacement; each interpolate joins one
$\Rnew$ draw and one pool draw.

\paragraph{Training of the classifiers.}
The \LRC{} and C2ST tests each train their own classifier with the class-balanced cross-entropy
of \cref{sec:method-competitors}: Adam with learning rate $10^{-3}$, moments $(0.9,0.999)$ and
weight decay $10^{-4}$, mini-batches of 128 windows per class, at most 1000 steps. Early
stopping holds out a random 20\% of each training half. This subset is drawn at random also
under the time-block split, so validation windows neighbour training windows in time and the
validation loss can be optimistic. The whole validation set is evaluated every 50 steps,
training stops after five evaluations without improvement, and the weights with the lowest
validation loss are restored. The selection affects power only, because the evaluation half is
never touched. In practice the classifiers are trained for very few steps. In E2 families
(a)--(d) the selected checkpoint was step 50 in all 4800 classifier fits at $n=500$ and in 4790
of 4800 at $n=2000$ (never later than step 100), and training stopped after a median of 300
steps. In the E3 pooled tests, with 1500 training windows per sample, the median selected step
was 100, with ranges 50--250 for C2ST and 50--200 for the \LRC{} over 25 fits each. Two features
of the procedure favour the first checkpoint. At $n=500$ the hold-out has 50 windows per class
(200 at $n=2000$), so the validation cross-entropy is a noisy criterion. And by step 50 the
classifier has drawn 50 mini-batches of 128 windows per class from 200 training windows per
class (800 at $n=2000$), so it can already fit the training windows closely, which is the
overfitting noted in the change log of the protocol. Since the first checkpoint is selected at
$n=2000$ as well, the noise of a small hold-out is not the only reason. The defaults were kept
after a pilot of 40 replications on alternatives outside the E2 grid (a GARCH persistence of
0.92 and Student-$t$ tails with $\nu=8$) that compared a learning rate of $10^{-4}$, a weight
decay of $10^{-2}$, training without early stopping and a minimum of 100 steps before a
checkpoint can be selected; according to the change log, none of these configurations had
higher power at valid size. The pilot is not part of the reported runs or of the repository.
Because the selected checkpoint is so early, E2 part f also fits the logit without early
stopping, for 1000 steps at learning rate $10^{-3}$ and for 3000 steps at $10^{-4}$
(\cref{tab:e2-llr}).

% ---------------------------------------------------------------------------------------------
\subsection{Time-block split and Driscoll--Kraay inference}\label{app:impl-split}

\paragraph{Assignment of blocks.}
The block labels of \cref{sec:method-inference} are computed from each sample's own range of
date ranks and live in one namespace shared by all samples of a call. The ten labels are put in
a uniformly random order, and following that order labels go to the training half until the
cumulative number of windows over all samples is closest to half of the total (30\% or 70\%
under the train-fraction settings of E5). If a half of some sample is empty, the order is
redrawn, at most 100 times. On the E3 data the 30 window dates of \TEST{} form ten blocks of
three dates, so every absolute test trains on five blocks and evaluates 1500 real and 1500
synthetic windows. The 62 window dates of $\Rpast$ form eight blocks of six dates and two of
seven, and the 123 window dates of an E4 \TEST{} sample form blocks of 12 or 13 dates. Training
and evaluation blocks are interleaved in random order, with no gap between adjacent blocks.

\paragraph{Grid and variance.}
The Driscoll--Kraay variance \eqref{eq:dk-var} is computed separately for the real and the
synthetic evaluation values, each on its own grid. A date rank $\tau$ is mapped to the grid
position $(\tau-\tau_{\min})/g$, where $\tau_{\min}$ is the first evaluation date of the sample
and $g$ is the greatest common divisor of the gaps between its distinct date ranks ($g=L$ for
non-overlapping windows). Grid positions without an evaluation window, in particular all
positions inside training blocks, enter the date sums $\mathcal S_u$ as zeros. Adding the two
variances in \eqref{eq:hac-z} presumes that the two samples are independent. This holds for a
synthetic against a real sample, and for $G$ against $\Rpast$ in the anchored test up to the
calibration caveat of \cref{sec:method-relative}. It fails for the ticker-half placebo of E1,
whose halves share every date: there the Driscoll--Kraay variance ignores the strong positive
same-date covariance of the halves, so those rows cannot reveal oversize, and the wild-bootstrap
rows, whose multipliers are shared by date, are the informative ones (\cref{sec:emp-placebo}).

\paragraph{Truncation lag.}
For each of the two series, $\hat\rho$ is the lag-one autocorrelation of the date sums over
adjacent grid positions that both hold evaluation windows, clipped to $[-0.97,0.97]$; at least
ten such pairs are required. The Andrews constant $\hat\eta$ and the lag $H$ are then those of
\cref{sec:method-inference}. With fewer than ten pairs the rule is applied to the sequence of
observed dates, and if that also fails the rule of \citet{newey1994automatic},
$H=\lfloor4(U'/100)^{2/9}\rfloor$ with $U'$ the number of evaluation dates, is used. The same
rule is the ``Newey--West'' option of E5; with 15 evaluation dates it gives $H=2$. In the E3
pooled \WC{} tests (five generators, five seeds) the Andrews rule chose $H$ between 0 and 3 for
the real and between 0 and 4 for the synthetic evaluation series.

\paragraph{Fixed-\texorpdfstring{$b$}{b} comparison rows (E1).}
The E1 tables add two fixed-$b$ p-values, computed after the runs from the stored $\zhac$, the
numbers of evaluation dates $U_X,U_Y$ of the two samples and their lags $H_X,H_Y$. Both refer
$\zhac$ to a simulated finite-sample version of the reference distribution of
\citet{kiefer_2005_new}: 40{,}000 draws of the HAC $t$-statistic, computed from \iid{} $N(0,1)$
series with the Bartlett kernel at the same lag, are symmetrised to 80{,}000 values, and the
p-value is $(1+\#\{\text{draws}\ge\zhac\})/(1+80{,}000)$. The \emph{conservative} version
simulates one series of length $\min(U_X,U_Y)$ with lag $\max(H_X,H_Y)$, which we expect, but
have not proved, to be the heaviest-tailed case of the two-sample statistic whatever the shares of
the two variances; a Monte Carlo check over the E1 designs supports it
(\cref{tab:e1-fixedb-check}). The
\emph{equal-variance} version simulates both samples with their own lengths and lags and equal
variances per date, which is the limit under an equal-law null with equal cross-sections, as
for the absolute and the anchored null of the simulated E1 placebo. Both treat the evaluation dates as a contiguous series, whereas the
Driscoll--Kraay estimator runs on a grid with zero-filled holes, so both are approximations.
Neither enters a verdict.

% ---------------------------------------------------------------------------------------------
\subsection{Kernel tests}\label{app:impl-kernel}

\paragraph{MMD.}
The kernel matrix of the pooled (subsampled) sample is built in row blocks with an exact
diagonal. The $B=500$ relabellings of the permutation test are processed in blocks of 100
through products of the kernel matrix with label-indicator matrices; ties with the observed
statistic count against it. In E4 the bandwidth
is fixed per universe at the median heuristic of the real \TEST{} windows, so that the five
generators of a universe are compared with one kernel.

\paragraph{Wild bootstrap.}
The multipliers of \cref{sec:method-inference} are generated as $\zeta_0\sim N(0,1)$ and
$\zeta_u=\phi\,\zeta_{u-1}+(1-\phi^2)^{1/2}e_u$ with $\phi=\exp(-1/\ell_w)$ and $e_u$ \iid{}
$N(0,1)$, on the grid of the pooled date ranks of both samples (built as above); $\ell_w=0$
gives \iid{} multipliers. The $B=500$ replicates are computed in blocks of 100. For the
correlation length $\ell_w$, at most 3000 pooled windows are drawn at random, and the five
leading eigenvectors of their doubly centred kernel matrix are computed by the Lanczos method.
The scores of all windows on these axes (by Nystr\"om extension for windows outside the subset)
are centred within each sample and summed by date, giving five series $a=1,\dots,5$. Let
$\hat\rho_a$ be the lag-one autocorrelation of series $a$ over adjacent observed dates (clipped
as above), $\hat\gamma_a$ its variance and $\hat s_a^2=\hat\gamma_a(1-\hat\rho_a^2)$. Then
\begin{equation*}
  \hat\eta=\frac{\sum_a4\hat\rho_a^2\hat s_a^4/\{(1-\hat\rho_a)^6(1+\hat\rho_a)^2\}}
                {\sum_a\hat s_a^4/(1-\hat\rho_a)^4},\qquad \ell_w=(\hat\eta\,U)^{1/3},
\end{equation*}
with $U$ the length of the pooled grid. This is the AR(1) plug-in constant of
\citet{andrews1991heteroskedasticity} with equal weights; for one series it reduces to the
$\hat\eta$ of \cref{sec:method-inference}. If the eigendecomposition fails, the multipliers are
\iid{}. In the E3 pooled tests $\ell_w$ ranged from 1.57 to 2.97 grid steps for MMD and from 1.01
to 1.82 for \SigMMD{}.

\paragraph{Signatures.}
A window is mapped to the $L+1$ points $(k/L,\sum_{l\le k}r_l)$, $k=0,\dots,L$, of a path in
$\R^2$. The depth-4 signature of a straight segment with increment $\Delta$ is the truncated
tensor exponential $(\Delta,\Delta^{\otimes2}/2!,\Delta^{\otimes3}/3!,\Delta^{\otimes4}/4!)$,
and the signature of the piecewise-linear path is the truncated tensor product of its $L$
segment signatures (Chen's identity; \citealp{lyons1998differential,chevyrev2016primer}),
computed for all windows at once. Levels 1 to 4 give $2+4+8+16=30$ features. The lead-lag
transform, available in the code, is not used. E4 runs only the \WC{} and MMD tests: the
depth-4 signature of the 31-dimensional time-augmented path of an E4 window would have 954{,}304
features, and the implementation refuses more than 50{,}000.

% ---------------------------------------------------------------------------------------------
\subsection{Model fitting}\label{app:impl-fit}

All optimisers are those of SciPy \citep{virtanen_2020_scipy}. Every fit uses the \TRAIN{}
returns of one ticker (E3, E5) or one universe (E4) in daily units.

\paragraph{GBM and Merton.}
GBM uses the sample mean and standard deviation (divisor $\Ttr$). Merton is fitted on returns
standardised by their sample mean and standard deviation, and the estimates are mapped back
exactly, since the model is location-scale equivariant. The mixture density is truncated at 11
jumps, the value of the truncation rule of \cref{sec:data-generators} at the upper intensity
bound 0.5, and evaluated with log-sum-exp. The likelihood is maximised by L-BFGS-B with the
analytic gradient (at most 1000 iterations, tolerances $10^{-12}$ on the objective and
$10^{-8}$ on the projected gradient) under the bounds $\lambda_J\in[10^{-4},0.5]$,
$\sigma\ge10^{-5}$ and $\sigma_J\ge10^{-4}$, with wide bounds on the standardised $\mu$ and
$\mu_J$. Six starting points are tried and the best optimum is kept: four moment-based starts
at $\lambda_J\in\{0.005,0.02,0.08,0.3\}$, in which the jump variance matches the excess
kurtosis, the jump mean the skewness and the diffusion variance the rest, a near-Gaussian start
($\lambda_J=10^{-3}$) and a start with frequent small jumps ($\lambda_J=0.4$). Parameters at a
bound are counted in \cref{tab:e3-fit-diagnostics}.

\paragraph{\texorpdfstring{GARCH-$t$}{GARCH-t}.}
The likelihood is maximised on returns divided by their sample standard deviation ($\mu$ and
$\omega$ are mapped back; $a$, $b$ and $\nu$ are scale invariant), with the in-sample filter
started at the sample variance. Under variance targeting \citep{francq_2011_merits} the free
parameters are $(\mu,a,b,\nu)$. They are estimated by SLSQP with the analytic gradient (at most
500 iterations, tolerance $10^{-11}$), subject to $a,b\ge0$, $a+b\le0.995$, $\nu\in[2.05,100]$
and $\mu$ within one standard deviation of the sample mean. The five starting values of
$(a,b,\nu)$ are $(0.05,0.90,8)$, $(0.10,0.85,5)$, $(0.15,0.70,6)$, $(0.03,0.96,10)$ and
$(0.08,0.60,4)$, and the best optimum is kept and rescaled onto the constraint $a+b\le0.995$ if rounding violates it. Out
of sample, the filter runs through \TRAIN{} and continues into \TEST{}. Simulations from the end
of \TRAIN{} start at the one-step variance forecast; stationary simulations start at the
unconditional variance and discard 250 days (1000 days for the E1 factor panel). Every GARCH
component in E1 and E3--E5 uses this estimator, targeting and cap included.

\paragraph{Heston.}
The variance follows the quadratic-exponential step of \citet{andersen_2008_simple} with a daily
time step and $\psi_c=1.5$, and the log-return the matching step with $\gamma_1=\gamma_2=1/2$
(trapezoidal weights for the integrated variance). The simulated method of moments averages
over $S=39$ paths of \TRAIN{} length, started from the stationary law of $v_t$; the code sets
$S=\min\{50,\max\{10,\lfloor50{,}000/\Ttr\rfloor\}\}$. Its random numbers are common to all evaluations of the objective
(fixed normal draws drive the variance step, its uniforms and the returns; fixed uniforms the
initial variance) and come from one fixed seed, the same for every ticker, so the objective is a
deterministic function of the parameters. The diagonal weights are the inverse variances of the
data moments over 200 stationary-bootstrap replicates with mean block length
$\max\{20,b_{\mathrm{PW}}\}$, where $b_{\mathrm{PW}}$ is the block-length rule of the bootstrap
below; the standard deviation of the tail quantile ratio is capped at 10\% of its value, which
floors its weight. The QE variance paths scale with $\theta_v$ and depend on $(\kappa_v,c)$
only, so the objective is profiled. For given $(\kappa_v,c)$, $(\log\theta_v,\rho_v)$ is
optimised by bounded Nelder--Mead (at most 400 evaluations), started from the variance-matching
$\theta_v$ and the best $\rho_v$ in $\{-0.8,-0.5,-0.2,0.1,0.4\}$. The profile in
$(\log\kappa_v,\log c)$ is evaluated on an $8\times8$ grid: eight values of $\kappa_v$
geometrically spaced from 0.003 to 0.4, and $c/\sqrt{1.5}$ equal to 0.15, 0.30, 0.45, 0.60,
0.72, 0.84, 0.93 and 1.00. It is then
minimised by bounded Nelder--Mead from the two best grid points (at most 120 profile evaluations
each) and once more from the better solution with a third of the initial steps. The estimate is
the best of all profile evaluations, including those at the starting grid points, and $\mu$ is set so that the
simulated mean return equals the sample mean. The starting variance of \TEST{} simulations is
the exponentially weighted average of squared \TRAIN{} returns with decay 0.94, started at the
sample variance. Fit times and parameters at bounds are in \cref{tab:e3-fit-diagnostics}, the
moment match in \cref{tab:e3-heston-moments}.

\paragraph{Stationary bootstrap.}
The mean block length follows the rule of \citet{politis2004automatic} with the correction of
\citet{patton2009correction}, implemented as in the \texttt{arch} package (version 9) and applied
to $\abs{r_t-\bar r}$, with two differences: the length is floored at one day, and 20 days are
used when the rule is degenerate, which never happened in E3 (\cref{tab:e3-fit-diagnostics}).

All univariate generators simulate the tickers independently: the path of ticker $i$ uses the
$i$-th child of the seed sequence of the synthetic seed. The synthetic panels of E3 and E5
therefore have no cross-sectional dependence, which the \texttt{indep} cell of E1 reproduces.

\paragraph{Multivariate generators (E4).}
The ECME iterations of MV Student-$t$ \citep{liu1995ml} update $\nu\in[2.1,200]$ by a bounded
Brent search over $\log\nu$ of the observed log-likelihood and stop when its relative increase
falls below $10^{-12}$ (at most 2000 iterations). CCC-GARCH-$t$ uses the \GARCHt{} estimator
above for every margin; with $F_{\nu_i}$ the unit-variance Student-$t$ distribution function of
margin $i$, the copula correlation $P$ is the sample correlation of the normal scores
$\Phi^{-1}(F_{\nu_i}(z_{it}))$ of the standardised residuals, and simulated innovations are
$F_{\nu_i}^{-1}(\Phi(x_{it}))$ with $x_t\sim N(0,P)$. Factor GARCH-$t$ applies the same
\GARCHt{} estimator to the factor and to each OLS residual. For the utility study, synthetic \TRAIN{}-length samples start from stationarity with 250
discarded days; for the tests and the placebo, \TEST{}-length samples start from the
end-of-\TRAIN{} state of every GARCH recursion, while the \iid{} models and the bootstrap are
stationary from the first day.

% ---------------------------------------------------------------------------------------------
\subsection{Seeds, checkpoints and software}\label{app:impl-seeds}

Each replication draws all its random numbers from a NumPy \texttt{SeedSequence}
\citep{harris_2020_array} determined by a fixed root seed per experiment and by the
replication's key (part, cell and replication index, hashed with SHA-256). Results therefore do
not depend on the number of worker processes, on how replications are grouped or on which parts
are run; in a check during development, E3's pooled tests and E2 family (a) run with different
numbers of workers gave identical outputs. Within a replication every test receives its own stream, derived from one
draw and the test's name, so a test's result does not depend on which other tests run beside
it. PyTorch is seeded from the replication's generator inside a forked random state. Synthetic
seeds are shared by all generators in E3 and E5, and the test stream of E4 is keyed by universe
and seed only, so generators are compared under common random numbers. The root seeds are
20251121 (E1), 20251115 (E2), 20251116 (E2-dep and E3), 20251118 (E4) and 20251117 (E5, which
reuses E3's synthetic seeds and power replications). The panel is drawn with seed 20251114,
which is also the fixed seed of the Heston calibration.

Replications are computed in blocks, and every finished block is written to disk, so that an
interrupted run resumes without recomputation. The checkpoint key hashes the configuration of
the part and, depending on the experiment, the library code (E1, E2, E5), the fitted parameters
(E3, E5) or the universes (E4), so a resumed run never mixes configurations. Each experiment
stores its configuration, a fingerprint of the code and the package versions: Python 3.12.16,
NumPy 2.5.4, SciPy 1.18.1, pandas 3.1.1, statsmodels 0.15.0 and PyTorch 2.15.0 (CPU build).
The \texttt{arch} package (version 9.0.0) serves only for cross-checks in the unit tests.

% ---------------------------------------------------------------------------------------------
\subsection{Compute and deviations from the planned design}\label{app:impl-compute}

The full runs of E1--E5 were executed on 15 November 2026 on one 16-core machine, shared with
other jobs and with about 6--7\,GB of free memory. Memory, not cores, was the binding
constraint: a worker process holds 0.46--0.9\,GB, the upper end for the kernel matrices of
$2\times3000$ windows. Each experiment process therefore ran three to five workers instead of
the twelve assumed in the protocol, with up to three processes at a time and one BLAS and
PyTorch thread per worker. According to the run logs, the replications took about 63 hours of
summed task time, the run time of each block of replications summed over blocks (E1 4.8, E2
27.7, E2-dep 2.1, E3 11.6, E4 9.6 and E5 7.2 hours), not
counting the model fits and the precomputed null distributions of the likelihood-ratio oracle,
and about 7.5 hours of wall-clock time from the start of E2 families (a)--(d) to the last table.
The architecture ablation ran on 16 November 2026 on the same machine, in about 24.3 hours of
summed task time and 3.1 hours of wall-clock time.

In the full runs of E3 the power check (designs D1 and D2) was completed before the first
real-data test started. The protocol was not fixed before every contact with the real data,
however. On the day of the full runs, development runs of the E3 script with six tickers, two
synthetic seeds, critics trained for 100 steps, 100 permutations and three replications per
power design exercised every part of the script, including the pooled real-data tests. Design
D2 and the combination of the seeds by $\pmed$ enter the change log of the protocol on that day,
and according to the file times of the development runs their first pooled tests preceded the
first run of D2. No output of the development runs is reported.

Before the runs, the replication counts of the experimental protocol (Section 4 of the design
document distributed with the code, and its change log) had been estimated at about 205 CPU
hours, against about 45 hours available on the day. The counts were therefore reduced, as
recorded in the change log of the protocol. \Cref{tab:impl-cuts} lists every deviation. Three of
them limit what the paper can claim. (i) Fewer seeds and replications widen every interval. For
example, the power of the decision rule in design D2, 0.45 from 100 replications, has the
Wilson interval $[0.36,0.55]$ (\cref{tab:e3-power}), and E3's verdict p-value is twice a median
over five rather than ten seeds. (ii) E2 part e shows the anchored
test at the boundary for $n=500$ only. (iii) Without the parts \texttt{gridsize},
\texttt{gridpower} and \texttt{gridrealvol}, E5 does not re-estimate the size of the absolute
test or the power and boundary size of the decision rule under any non-default configuration.
A change of an absolute decision in the E5 grid can then not be separated from a change of size,
and because the power is unknown, a non-rejection in the grid is \INCONCL{} by construction: no
grid configuration can return \GOOD{}. The size of the default real-data pipeline is measured by
E1.

\begin{table}[tbp]
  \centering
  \caption{Planned and executed replication counts. ``Planned'' is Section 4 of the design
    protocol or its change log or, where the protocol gives no number, the default of the
    experiment script. ``As planned'' means no reduction. Seeds are synthetic seeds (independent
    synthetic samples and test randomness); each seed of the E3 and E5 power checks is a full
    run of the anchored test.}
  \label{tab:impl-cuts}
  \footnotesize
  \begin{tabular}{@{}>{\raggedright\arraybackslash}p{0.30\linewidth}
                  >{\raggedright\arraybackslash}p{0.33\linewidth}
                  >{\raggedright\arraybackslash}p{0.31\linewidth}@{}}
    \toprule
    Experiment, part & Planned & Executed \\
    \midrule
    E1 (a), simulated panel placebo & 200 replications & 200 per cell (cells \texttt{joint},
      \texttt{indep}) \\
    E1 (b), ticker-half placebo & 100 splits in \TEST{} and 100 in \TRAIN{} & 50 and 50 \\
    E1 (c), \TEST{} vs \TRAIN{} & 10 test seeds (script) & as planned \\
    \addlinespace
    E2 (a)--(d) & $n\in\{500,2000\}$; 200 replications under $H_0$, 100 per alternative &
      as planned \\
    E2 (e), relative test at the boundary & $n\in\{500,2000\}$, 200 replications &
      $n=500$ only, 200 replications \\
    E2 (f), learned vs exact log-likelihood ratio & 20 replications, three training budgets
      (script) & as planned \\
    E2 (g), dependent windows & $n=500$, 200 replications & as planned \\
    E2-dep (added after the first E2 results) & $n\in\{500,2000\}$; 200 replications under
      $H_0$, 100 per alternative (script) & $n=2000$ only; 100 and 50 \\
    \addlinespace
    E3, pooled tests & 10 seeds & 5 seeds; $\pmed$ over 5 seeds \\
    E3, power check D1 and D2 & 200 replications per design (minimum 50), 10 seeds each &
      100 per design, 5 seeds each \\
    E3, per-ticker null & 10 pairs per ticker (script) & as planned \\
    \addlinespace
    E4, tests & 10 seeds & 2 seeds \\
    E4, placebo & 200 replications per generator & 60 per generator \\
    \addlinespace
    E5, one-at-a-time grid & 14 configurations, 10 seeds & 14 configurations, 5 seeds \\
    E5, alternative split & 10 seeds; rule power 50 replications per design & 5 seeds; 25 per
      design \\
    E5, critic training on simulated data & 100 replications per cell & 50 per cell \\
    E5, size and rule power per configuration (\texttt{gridsize}, \texttt{gridpower},
      \texttt{gridrealvol}) & 200 replications per size design; 50 per power design (100 for
      the default) & not run \\
    \addlinespace
    Workers per process & 12 & 3--5 (memory) \\
    \bottomrule
  \end{tabular}
\end{table}

% ---------------------------------------------------------------------------------------------
\subsection{Availability and provenance}\label{app:impl-avail}

\paragraph{Code and records.}
The package \texttt{synthval}, the experiment scripts, the command list
\texttt{experiments/run\_full.sh}, the result files under \texttt{results/} and the scripts in
\texttt{paper/scripts} that write the condensed tables of this paper are kept in one git
repository. Next to its results, each experiment writes meta files
(\texttt{results/<experiment>/meta*.json}) that record the configuration and root seed, the
command line, the number of workers, the git commit and the list of uncommitted files at run
time, a SHA-256 fingerprint of the source files of \texttt{synthval}, the package versions of
\cref{app:impl-seeds} and the SHA-256 of the price file. The run logs, the checkpoints and the
development runs are kept outside version control; the compute figures of
\cref{app:impl-compute} are taken from the logs.

\paragraph{Code state of each result.}
Every part of E1--E5 ran with the same library fingerprint (prefix \texttt{32ae8a76}). All parts
but one record commit \texttt{b7c16b8}, with at most two files of the experiments directory not
committed: the command list (modified) and the script of E2-dep (new). Both are committed in the
next commit, \texttt{38595f6}: the meta files of the ablation, run at that commit, no longer list
them as uncommitted.
The utility part of E4 records the earlier commit \texttt{2cb296e}, at which the library and the
experiment scripts were not yet committed, with the same library fingerprint. The top-level meta
file of each experiment, written by the final pass that produces its tables, records the SHA-256
of the experiment script, and for all six scripts it equals that of the script in the
repository. The command list arranges the runs in three lanes. E2 families (a)--(d) were started
earlier in a separate shell, and E2-dep and the real-data parts of E1 ran in a fourth shell while
lane A was busy, so lane A found their results in the checkpoints. The architecture ablation ran
at commit \texttt{38595f6} with uncommitted changes to the network module and, for its simulated
part, to the critic and classifier modules, and with its script not yet committed. Its library
fingerprints (prefixes \texttt{40704594} for the real-data part and \texttt{ada27d81} for the
simulated part) therefore differ from that of E1--E5; the changes were committed afterwards
(commit \texttt{764c44f}) together with a device option, and the script in the repository has
the SHA-256 recorded in the ablation's meta file.

\paragraph{Device.}
All reported runs, the ablation included, used the CPU build of PyTorch 2.15.0 with one thread
per worker. The device option, which no reported run uses, moves the networks to a GPU
(environment variable \texttt{SYNTHVAL\_DEVICE}, default CPU) while every random draw (initial
weights, mini-batch indices, interpolation weights) stays on the CPU generator. A GPU run then
sees the same mini-batches and initial weights, but differs by floating-point rounding and by
non-deterministic GPU kernels, so it need not reproduce the reported numbers digit for digit.

\paragraph{Data.}
The prices are a Bloomberg export: the fields \texttt{PX\_LAST} and \texttt{PX\_VOLUME} of the
503 tickers that were S\&P~500 constituents at extraction, requested from 2018-01-01 to
2025-06-17. Bloomberg data are licensed, so the price file is not distributed
with the code. A user with Bloomberg access can rebuild the inputs from the 100 tickers of
\cref{tab:data-panel} (E1, whose factor panel is fitted to them, E3 and E5) and from the E4
universes, listed in \texttt{results/e4/universes.csv}, with the rules of \cref{app:data}. The
SHA-256 recorded in every meta file,
\texttt{ddfd0d23dc274a58\allowbreak f8089e53b3872754\allowbreak 005033447d77eb26\allowbreak
50a0795382f832dc}, identifies the extract that was used. E2 and E2-dep use no data.

\paragraph{Pilots.}
Three choices rest on checks that are recorded only in the change log of the protocol and cannot
be reproduced from the repository: the learning rate $10^{-3}$ of the critic (an informal check
before any E2 run, \cref{sec:method-critic}), the defaults of the classifiers (the pilot of 40
replications of \cref{app:impl-nets}) and the time-pooled architecture of the ablation (a pilot
on four E2 designs at $n=500$). The first two fixed settings of the pre-specified tests; the
third concerns only the post-hoc ablation. The development runs of the experiment scripts, with
reduced settings, are also kept outside version control (\cref{app:impl-compute}).
